\section{Conclusion}
\label{sec:conclusion}

Historical assets are usually valued for remembering what a system already
knew how to represent. This paper proposes a stronger and more demanding role:
a governed causal Episode corpus can preserve the option to discover that the
old representation itself was inadequate.

That option requires structure and restraint. KFD-4 makes replay
perspective-bound and loss-aware. KFD-5 keeps candidate genesis open while
making qualification answerable to facts, alternatives, deletion, falsifiers,
and real work. Draft KFD-6 adds plural method comparison, fixed-ontology and
no-new-Primitive baselines, held-out evaluation, false-candidate retention,
and separated promotion authority. Episode storage alone performs none of
these acts.

The current KFD case registry makes the research direction more visible. It
contains six real candidate tracks, including three bounded software-domain
acceptances and a retained natural pre-enactment of agent-generated ontology
change. It also makes the missing evidence visible: no preregistered plural
method experiment, no systematic false-candidate result, no held-out transfer,
and no autonomous ontology-review trigger. That is why the evidence supports
feasibility, not completion.

The practical consequence is significant. If agents can participate in
defining the objects through which work is understood, problem definition
becomes an operational knowledge-production surface. Preserving that surface
can expand who notices unnamed burdens; governing it poorly can concentrate
power in corpus, model, or promotion owners. The right response is neither
unbounded retention nor unaccountable automation. It is a verifiable
separation among causal history, perspective, candidate generation,
qualification, independent evaluation, and admission.

The next milestone is therefore empirical: run the preregistered ablation and
shadow-loop protocols in Section~\ref{sec:evaluation}, retain the negative
results, and let failure narrow the claim. If a candidate survives those gates
and improves held-out continuation without hiding participant or source
authority, Episode history will have demonstrated value not merely as memory,
but as infrastructure for revising what a system is capable of knowing.
